% Pollis drill-down: Prediction Competitions, the package PollisCompetitions.jl, its TOML files, its DuckDB database and
% the folders that Manage Competition Datasets, Score Models and Agents' Forecasting Competition write (20 November 2026).
% For users who want the detail, and for whoever maintains Pollis.
% Compile with: latexmk -pdf workflow_prediction_competitions.tex
% Shared look for the one-page Pollis workflow handouts (same fonts and colours as the brochure).
\documentclass[10pt,a4paper]{article}

\usepackage[a4paper,left=17mm,right=17mm,top=14mm,bottom=16mm,footskip=8mm]{geometry}
\usepackage[T1]{fontenc}
\usepackage[default,semibold]{sourcesanspro}
\usepackage{microtype}
\usepackage{graphicx}
\usepackage[table]{xcolor}
\usepackage{tikz}
\usepackage[most]{tcolorbox}
\usepackage{fontawesome5}
\usepackage{enumitem}
\usepackage{fancyhdr}
\usepackage[hidelinks]{hyperref}

\graphicspath{{images/}}

% ---------------------------------------------------------------- colours (as in the brochure)
\definecolor{navy}{HTML}{1C3A5E}
\definecolor{pblue}{HTML}{3B78B5}
\definecolor{amber}{HTML}{F2B33D}
\definecolor{deepamber}{HTML}{D9902A}
\definecolor{pale}{HTML}{EEF3F9}
\definecolor{ink}{HTML}{2A2F36}
\definecolor{logonavy}{HTML}{173D63}
\definecolor{logodark}{HTML}{C98A1E}
\definecolor{logoamber}{HTML}{F2B446}
\definecolor{logocream}{HTML}{FEEFD2}
\definecolor{soft}{HTML}{5E6773}
\color{ink}

\setlength{\parindent}{0pt}
\setlength{\parskip}{3pt}
\setlist{nosep,leftmargin=*,itemsep=2pt}

\pagestyle{fancy}
\fancyhf{}
\renewcommand{\headrulewidth}{0pt}
\fancyfoot[L]{\footnotesize\color{soft}\textbf{Pollis}\enspace Bright Insights through Exploration, Modelling, Simulation and Optimisation}
\fancyfoot[R]{\footnotesize\color{soft}github.com/saragga/pollis/releases/latest}

% The Pollis logo, redrawn so it prints sharply at any size
\newcommand{\pollislogo}[1]{%
  \begin{tikzpicture}[x=#1/512,y=-#1/512,baseline=(current bounding box.center)]
    \fill[logonavy] (0,0) rectangle (512,512);
    \fill[logodark] (36,422) -- (104,256) -- (172,422) -- cycle;
    \fill[logoamber] (104,422) -- (203,180) -- (302,422) -- cycle;
    \fill[logocream] (203,422) -- (339,89) -- (476,422) -- cycle;
  \end{tikzpicture}}

% Top of the page: logo, name, kind of document (Workflow unless given, e.g. [Drill-Down]) and its subject;
% then the question and a one-line promise
\newcommand{\workflowhead}[4][Workflow]{%
  \thispagestyle{fancy}%
  \begin{minipage}[c]{13mm}\pollislogo{11mm}\end{minipage}%
  \begin{minipage}[c]{\dimexpr\linewidth-13mm\relax}
    {\color{navy}\Large\bfseries Pollis}\enspace{\color{soft}#1}\hfill
    {\color{deepamber}\small\bfseries\MakeUppercase{#2}}
  \end{minipage}\par\vspace{7mm}
  {\color{navy}\fontsize{22}{26}\selectfont\bfseries #3}\par\vspace{4pt}
  {\color{soft}\large #4}\par\vspace{5mm}}

% A menu path: \menu{Explore}{Access Data Libraries}... is written as \path{Explore \gt Access Data Libraries}
\newcommand{\gt}{\,{\color{deepamber}\bfseries ›}\,}
\newcommand{\menupath}[1]{{\footnotesize\bfseries\color{navy}#1}}

% A screenshot if images/<file> exists, otherwise an orange placeholder that says what to capture
\newlength{\shotheight}\setlength{\shotheight}{43mm}
\newcommand{\stepshot}[2]{%
  \IfFileExists{images/#1}{%
    \tcbox[enhanced, colback=black, colframe=black!70, boxrule=0.4pt, arc=2pt,
      left=0pt,right=0pt,top=0pt,bottom=0pt, boxsep=0pt, halign=center]{%
      \includegraphics[width=\dimexpr\linewidth-1pt\relax,height=\shotheight,keepaspectratio]{#1}}}{%
    \begin{tcolorbox}[enhanced, colback=white, colframe=orange!80, boxrule=0.6pt, arc=2pt,
      height=\shotheight, valign=center, halign=center, left=6pt, right=6pt]
      {\color{orange}\Large\faCamera}\par\smallskip
      {\color{orange}\footnotesize\bfseries Screenshot: #2}\par\smallskip
      {\color{orange!80!black}\scriptsize images/#1}
    \end{tcolorbox}}}

% One step: number, title, menu path, image file, what to capture, what it means
\newcommand{\step}[6]{%
  \begin{tcolorbox}
    {\color{deepamber}\Large\bfseries #1}\enspace{\color{navy}\large\bfseries #2}\par\vspace{1pt}
    \menupath{#3}\par\vspace{4pt}
    \stepshot{#4}{#5}\par\vspace{3pt}
    {\small #6}
  \end{tcolorbox}}

\newenvironment{steps}{%
  \begin{tcbraster}[raster columns=2, raster equal height=rows, raster column skip=4mm,
    raster row skip=4mm, enhanced, colback=pale, colframe=pale, boxrule=0pt, sharp corners,
    left=6pt,right=6pt,top=6pt,bottom=6pt, fontupper=\raggedright]}{\end{tcbraster}}

% Boxes for the bottom of the page
\newtcolorbox{pbox}[2][]{enhanced, colback=pale, colframe=amber, boxrule=0pt,
  leftrule=3pt, sharp corners, left=7pt, right=7pt, top=6pt, bottom=6pt,
  fontupper=\small\raggedright, before upper={{\color{navy}\normalsize\bfseries #2}\par\smallskip}, #1}
\newtcolorbox{nbox}[1][]{enhanced, colback=navy, colframe=navy, boxrule=0pt,
  sharp corners, left=8pt, right=8pt, top=6pt, bottom=6pt, coltext=white,
  fontupper=\small\raggedright, #1}

% A Next Steps item inside the bottom boxes
\newcommand{\nsitem}[2]{\item[\color{deepamber}\faChevronRight] \textbf{#1}: #2}

% Citations: \dcite{key}{Breiman (2001)} links to its entry \dref{key}{...} in the References, keyed as in pollis.bib
\newcommand{\dcite}[2]{\hyperlink{ref:#1}{\color{pblue}#2}}
\newcommand{\dref}[2]{\raisebox{1.1em}[0pt][0pt]{\hypertarget{ref:#1}{}}#2\par}
\newcommand{\ddoi}[1]{\href{https://doi.org/#1}{\color{pblue}doi:#1}}
\newcommand{\references}{\section*{References}\addcontentsline{toc}{section}{References}\footnotesize\raggedright\setlength{\parskip}{3pt}}

% The page of each heading, for Ask Pollis
% Pollis: the page of each heading, for Ask Pollis. A document that inputs this file, at the end of its preamble, writes
% <jobname>.labels beside its PDF, a line page|level|heading for each heading, the page counted from the first, the
% cover, as 1. compile_documents.jl copies it to ~/.pollis/documents, and ask_index.jl cuts the document into passages
% under its headings with it. \asklabel{1}{Title} marks a topic or a page, \asklabel{2}{Title} a heading inside it; the
% sections and subsections of a document that has them are marked as they go into its contents. A package vignette
% inputs it before its \documentclass (compile_vignette.jl), so the contents are taken over when the document begins,
% after the packages have taken them over. Read a second time, as by a copy of a document for Edit the LaTeX Source,
% which inputs it as resources/ask_labels when Pollis has put it in already (document_source.jl), it does nothing.
\ifdefined\asklabel\endinput\fi
\makeatletter
\newwrite\pollis@labels
\immediate\openout\pollis@labels=\jobname.labels
\newcommand{\asklabel}[2]{\write\pollis@labels{\the\ReadonlyShipoutCounter|#1|\unexpanded{#2}}}
\def\pollis@toc{toc}
\newcommand{\pollis@contentsline}[3]{%
  \def\pollis@file{#1}%
  \ifx\pollis@file\pollis@toc\asklabel{#2}{#3}\fi
  \pollis@addcontentsline{#1}{#2}{#3}}
% The headings of the New Methods style of the vignettes, \nmsection and \nmsub, which go into no contents
\def\pollis@asksection#1{\pollis@nmsection{#1}\asklabel{1}{#1}}
\def\pollis@asksub#1{\pollis@nmsub{#1}\asklabel{2}{#1}}
\AddToHook{begindocument/end}{%
  \let\pollis@addcontentsline\addcontentsline\let\addcontentsline\pollis@contentsline
  \ifdefined\nmsection\let\pollis@nmsection\nmsection\let\nmsection\pollis@asksection\fi
  \ifdefined\nmsub\let\pollis@nmsub\nmsub\let\nmsub\pollis@asksub\fi}
\makeatother



\usepackage{titlesec}
\usepackage{tabularx}
\usepackage{booktabs}
\usepackage{listings}

\fancyhf{}
\fancyfoot[L]{\footnotesize\color{soft}\textbf{Pollis}\enspace Drill-down: Prediction Competitions}
\fancyfoot[R]{\footnotesize\color{soft}\thepage}

\titleformat{\section}{\color{navy}\large\bfseries}{\color{deepamber}\thesection}{0.6em}{}
\titleformat{\subsection}{\color{navy}\normalsize\bfseries}{\color{deepamber}\thesubsection}{0.5em}{}
\titlespacing*{\section}{0pt}{12pt}{4pt}
\titlespacing*{\subsection}{0pt}{8pt}{2pt}
\setlength{\parskip}{4pt}

\lstset{basicstyle=\ttfamily\footnotesize, columns=fullflexible, keepspaces=true, breaklines=true,
  backgroundcolor=\color{pale}, frame=leftline, rulecolor=\color{amber}, framerule=2pt, xleftmargin=6pt,
  framexleftmargin=6pt, aboveskip=4pt, belowskip=4pt, upquote=true, tabsize=2}

\usepackage{xurl}
\newcommand{\file}[1]{{\small\nolinkurl{#1}}}
\newcommand{\sig}[1]{{\small\ttfamily\detokenize{#1}}}
\newcolumntype{L}{>{\raggedright\arraybackslash}X}

\begin{document}

\workflowhead[Drill-Down]{Prediction Competitions}
  {The Test Set Decides}
  {A forecast or a prediction is judged on values its maker never saw. This is how Pollis keeps them apart: the
   package PollisCompetitions.jl and its parts, the TOML files that describe datasets, models and scores, the DuckDB
   database built from them, and the folders written on your computer and in your Pollis Workspace.}

\begin{nbox}
{\color{amber}\bfseries Drill-down.}\enspace
A look under the surface of the menu \textbf{Compose\gt Prediction Competitions}: Manage Competition Datasets, Score
Models and Agents' Forecasting Competition. It is for anyone who wants to know where each file is and what it holds,
and for whoever maintains Pollis. It describes the state of 20 November 2026.
\end{nbox}

\tableofcontents

\section{Prediction Competitions in one table}

Three pages share one package and one rule. Every dataset is found, read and scored the same way, whoever contributes
it: the package, a toolbox or you.

\begin{center}\footnotesize
\renewcommand{\arraystretch}{1.25}
\begin{tabularx}{\linewidth}{@{}p{0.13\linewidth}LLL@{}}
\toprule
& \textbf{Manage Competition Datasets} & \textbf{Score Models} & \textbf{Agents' Forecasting Competition}\\
\midrule
Does & Lists the Dataset Catalog; you tick the datasets that are Active, save one to the workspace, or add one of your own
  & Runs a Julia function \sig{forecast(y, h, period)} of yours on the series of the active datasets and scores it
  & Several coding agents each write a forecasting method for the same series; Pollis scores them, and their mean and median\\
Reads & \file{metadata/datasets.toml}, the toolboxes' datasets, each \file{dataset.toml} of the workspace, the schema \file{forecastingCatalogue.toml}
  & The \file{info.toml} and Parquet files of each dataset; \file{run.toml}, \file{progress.json}, \file{results.json}
  & The same data; \file{study.toml}, \file{prepared.toml}, each agent's \file{method.jl} and \file{REPORT.md}\\
Writes & \file{data/forecasting/} of the Pollis Workspace
  & \file{~/.pollis/competitions/runs/<run>/}
  & \file{orchestration/forecasting-competition/<study>/} of the Pollis Workspace\\
Scores with & \textendash & \file{PC.score}, \file{PC.benchmarks}, \file{PC.step_errors} & \file{PC.score}, \file{PC.benchmarks}\\
\bottomrule
\end{tabularx}
\end{center}

Pollis keeps one setting of its own: the datasets ticked Active, in the profile's storage
(\file{pollis.manageCompetitionDatasets.active}), which Score Models lists in its step 1.

\section{The package PollisCompetitions.jl}
\label{sec:package}

PollisCompetitions.jl is a Julia package of PollisRegistry, installed with \file{Pkg.add("PollisCompetitions")} the
first time a page needs it. It holds three catalogues as TOML, the data of the built-in datasets as Parquet, the
functions that read and score them, and a module of eight statistical benchmarks. DuckDB
\dcite{RM2019a}{(Raasveldt and Mühleisen, 2019)} does the reading and the arithmetic: the Parquet files are read where
they lie, never copied into a database file.

\subsection{Its folders}

\begin{lstlisting}
PollisCompetitions.jl/
  Project.toml
  README.md
  src/
    PollisCompetitions.jl      the module: includes the files below, in this order
    benchmarks/
      Benchmarks.jl            the module Benchmarks
      ets.jl                   simple exponential smoothing and Holt's method, fitted by Nelder-Mead
      methods.jl               the seasonality test, the seasonal index, Theta and the eight benchmarks
    monash.jl                  read_tsf: the .tsf format of the Monash archive
    database.jl                the schema, connect_db, describe_db, load_table, load_dataset
    parquet.jl                 writing tables to Parquet through DuckDB
    data.jl                    where a dataset is: folder, info, competition
    tables.jl                  reading a dataset: series, observations, tables, table
    measures.jl                sMAPE, MASE and OWA: score, step_errors
    dataset_benchmarks.jl      the benchmarks on the series of a dataset
  metadata/
    datasets.toml              the Dataset Catalog: 132 datasets
    models.toml                the Model Catalog: 268 models
    benchmarks.toml            the Benchmarks: 5,828 runs
  data/builtin/<dataset_id>/   the 60 built-in datasets: 10 of series, 50 tabular
  test/runtests.jl
\end{lstlisting}

\subsection{Its functions}

The package is used as \sig{import PollisCompetitions as PC}. The functions below are those the pages and your code
call; the others are the parts they are built from.

\begin{center}\footnotesize
\renewcommand{\arraystretch}{1.2}
\begin{tabularx}{\linewidth}{@{}p{0.36\linewidth}L@{}}
\toprule
\textbf{Function} & \textbf{What it gives}\\
\midrule
\multicolumn{2}{@{}l}{\textit{Where a dataset is (data.jl)}}\\
\file{folder(id)} & \file{data/builtin/<id>} of the package if it exists, else \file{~/.pollis/competitions/<id>}; for the folder of a dataset of your own, that folder\\
\file{info(id)} & Its \file{info.toml}, or \file{nothing} when its data are not on this computer\\
\file{competition(id)} & Its id, name and format, the name of the Dataset Catalog when \file{info.toml} gives none\\
\file{is_dataset(path)} & Whether \file{path} is a folder holding a dataset of your own (\file{series.parquet} and \file{observations.parquet})\\
\multicolumn{2}{@{}l}{\textit{Reading a dataset (tables.jl)}}\\
\sig{series(id; frequency, domain)} & The series: number, name, frequency, period, horizon, domain, start\\
\sig{observations(id; series, frequency, split)} & The values: series, t, value and split, \file{train} or \file{test}\\
\file{tables(id)}, \sig{table(id, name; filters...)} & The Parquet files of a dataset, and one of them filtered by its columns, such as \sig{table("heloc", "splits"; repeat = 1, fold = 2)}\\
\sig{query(sql, parameters)} & Any SQL, run by DuckDB in memory\\
\multicolumn{2}{@{}l}{\textit{Scoring (measures.jl, dataset\_benchmarks.jl)}}\\
\sig{score(id, forecasts)} & sMAPE, MASE and OWA of a table of series, h and value, by frequency, by domain when the dataset gives domains, and in total\\
\sig{step_errors(id, forecasts)} & The sMAPE and the MASE of each series at each step ahead\\
\sig{benchmarks(id; methods, series)} & The forecasts of the benchmarks named, from \file{BENCHMARK_METHODS}, on the series given or all\\
\multicolumn{2}{@{}l}{\textit{The database (database.jl)}}\\
\file{connect_db(path)} & The database of \S\ref{sec:database}, in memory or in the file \file{path}\\
\file{load_table(name)} & One of the five tables, or the view \file{model_comparisons}\\
\file{load_dataset(id)} & The \file{data.parquet} of a built-in tabular dataset\\
\file{read_tsf(path)} & A \file{.tsf} file of the Monash archive, read into series and values\\
\bottomrule
\end{tabularx}
\end{center}

\subsection{The module Benchmarks}

\file{PollisCompetitions.Benchmarks} forecasts a single series from its values alone, and never reads a dataset. That
is why the agents of Agents' Forecasting Competition may use it and must not use the rest of the package. It exports
\file{BENCHMARKS}, \file{benchmarks}, \file{holt}, \file{is_seasonal}, \file{seasonal_index}, \file{ses} and
\file{theta}.

\begin{itemize}
  \item \textbf{Naive} repeats the last value; \textbf{sNaive} the last period of values.
  \item A series that passes \file{is_seasonal} is divided by its classical multiplicative seasonal index
    (\file{seasonal_index}); \textbf{Naive2}, \textbf{SES}, \textbf{Holt}, \textbf{Damped} and \textbf{Theta}
    \dcite{AN2000a}{(Assimakopoulos and Nikolopoulos, 2000)} forecast the adjusted values, which are multiplied back by
    the last period of the index.
  \item \textbf{Com} is the mean of SES, Holt and Damped.
\end{itemize}

\sig{benchmarks(y; period, horizon, methods)} gives the forecasts of each method named. A method that cannot be fitted
to a series raises an error, and the pages then leave that benchmark out of the table rather than score it on fewer
series.

\subsection{How a forecast is scored}

\file{score} works in a DuckDB database in memory, in four steps.
\begin{enumerate}
  \item \file{prepare!} makes three tables: \file{known_series}; \file{test_values}, the held-out values with \file{h}
    from 1; and \file{scales}, the mean absolute difference between each training value and the one a seasonal period
    before, the denominator of MASE \dcite{HK2006a}{(Hyndman and Koehler, 2006)}.
  \item \file{given!} reads your forecasts and refuses series the dataset does not have.
  \item \file{errors!} gives the sMAPE and the MASE of each series, the means over its horizon, and refuses a series
    not forecast once at each step.
  \item \file{reference!} computes Naive2 on the same series. \file{summary} averages the errors by frequency, by
    domain and in total, and weighs them against Naive2 into OWA
    \dcite{MSA2020a}{(Makridakis, Spiliotis and Assimakopoulos, 2020)}.
\end{enumerate}

The result is a table of \file{measure}, \file{breakdown}, \file{level} and \file{value}, where \file{level} is a
frequency, a domain or \file{Total}.

\section{The TOML files}
\label{sec:toml}

Every file is plain TOML and edited by hand. A field left out is unknown, and loads as NULL (\file{missing} in Julia).

\subsection{The Dataset Catalog: metadata/datasets.toml}

There is one \file{[[datasets]]} table for each dataset, with its works as \file{[[datasets.references]]}.

\begin{lstlisting}
[[datasets]]
dataset_id = "tourism_2011"                 # the key, and the name of its folder
dataset_name = "Tourism Forecasting Competition"
competition_name = "Tourism Forecasting Competition"   # given with competition_year, or neither
competition_year = 2011
topic = "Tourism"
data_structure = "Time-Series"              # Cross-section | Time-Series | Panel
geometry = "None"                           # None | Spatial Weights Matrix | Coordinates | Polygons
response_type = "Regression"                # Binary | Multiclass | Ordered | Count | Survival | Regression
rows = 164581
features = 0
size_mb = 0.45                              # Parquet files if built in; compressed download if not
licence = "CC BY 4.0 (via the Monash archive on Zenodo)"
delivery = "Built-in"                       # Built-in | Download
source_url = "https://..."

  [[datasets.references]]
  bibtex_key = "AHSW2011a"                  # a key of PollisReferences
  role = "Competition"                      # Primary | Competition | Benchmark | Other
\end{lstlisting}

The same \file{[[datasets]]} tables come from two other contributors. An installed toolbox names a file of them in
\file{forecastingDatasets} of its \file{pollisToolboxes} contribution, and a dataset you add has them in its
\file{dataset.toml} (\S\ref{sec:workspace}). Manage Competition Datasets reads all three and checks each against the
schema that Pollis ships, \file{compose/browser/media/forecastingCatalogue.toml}. That schema has 17
\file{[[field]]} tables (name, label, group, type, \file{controlled} and \file{values}, description) and four
\file{[[view]]} tables, the column presets of the page (id, label, \file{pinned}, \file{columns}). The page's own TOML
adds the curated collections as \file{[[collections]]} (id, name, year, url, description, \file{datasets} and
\file{papers}).

\subsection{The Model Catalog: metadata/models.toml}

\begin{lstlisting}
[[models]]
model_id = "tabarena_limix_2_1016_default"
model_name = "LimiX-2-1016 (default)"
estimator_family = "Foundation Model"       # one of nine families
estimator_method = "Pretrained in-context learner"   # one of the methods of its family
package = "LimiX"
language = "Python"                         # Julia | Python | R | C/C++ | Fortran

  [[models.references]]
  bibtex_key = "ZRYZT...2026a"
  role = "Methodology"                      # Methodology | Software | Other
\end{lstlisting}

There are nine families: Foundation Model, Tree-based, Neural Network, Baseline, Statistical, Kernel Methods, Fuzzy
System, Hybrid and Other. Each has a closed list of methods, which the header of \file{models.toml} gives and
\file{ESTIMATOR_METHODS} checks. A Hybrid model's method names the methods it combines.

\subsection{The Benchmarks: metadata/benchmarks.toml}

Each run of a model on a dataset is a \file{[[runs]]} table, with its scores under \file{[runs.scores]}, one
\sig{metric = score} line each.

\begin{lstlisting}
[[runs]]
run_id = "tabarena_catboost_default.airfoil_noise"   # <model_id>.<dataset_id>[.<subset>]
model_id = "tabarena_catboost_default"
dataset_id = "airfoil_noise"
split_protocol = "TabArena: 10 repeats of 3-fold cross-validation"

  [runs.scores]
  rmse = 1.5782966827503813
  time_train_s = 49.54921263853709
  time_infer_s = 0.04112234115600586
\end{lstlisting}

Comment blocks in the file say where each group of runs comes from: which publication, which splits and which metric.

\subsection{The facts of a dataset: info.toml}

Each dataset's folder has an \file{info.toml}, written by the script that converted it. Its \file{format} tells the
package how to read the folder.

\textbf{A dataset of series}, \sig{format = "monash"}, the format of the Monash Time Series Forecasting Archive
\dcite{GBWHM2021a}{(Godahewa, Bergmeir, Webb, Hyndman and Montero-Manso, 2021)}:

\begin{lstlisting}
competition = "tourism_2011"
name = "Tourism Forecasting Competition"
format = "monash"
converted = "2026-11-19T23:54:20.244"
fingerprint = "1be847ba..."                 # SHA-256 of the source files

[tables]                                     # the rows of each Parquet file
observations = 164581
series = 1311

[[frequencies]]                              # one for each frequency, in the order of the series
name = "Yearly"                              # left out for series that have no frequency
series = 518
shortest_horizon = 4
longest_horizon = 4
\end{lstlisting}

\textbf{A tabular dataset}, \sig{format = "tabular"}, with its response and the cross-validation it is scored by:

\begin{lstlisting}
competition = "heloc"
name = "heloc"
format = "tabular"
target = "RiskPerformance"                   # the response column of data.parquet
openml_task = 363676
openml_dataset = 46932
repeats = 10
folds = 3
tabarena_repeats = 3
converted = "2026-11-19T20:06:18.846"
fingerprint = "4cadc94a..."

[tables]
data = 10459
splits = 104590
\end{lstlisting}

Score Models lists a dataset, and Manage Competition Datasets shows it as Ready in its column Model Scoring, when its
\file{info.toml} has \sig{format = "monash"} and at least one \file{[[frequencies]]} table. The step Datasets shows
the horizon of a frequency as one number, or as a range when \file{shortest_horizon} and \file{longest_horizon} differ.

\section{The database}
\label{sec:database}

\file{PC.connect_db()} builds the database in memory, or in a file when given a path. PollisDB opens it as
\file{PollisCompetitions} in the Databases view. It has the five metadata tables, loaded from the three TOML files;
the view \file{model_comparisons}; and a schema for each dataset whose data are on this computer, with a view for each
of its Parquet files. Every table, view and column carries a comment (\file{describe_db}), which the Databases view
shows.

\subsection{The five tables and the view}

\begin{lstlisting}[language=SQL]
CREATE TABLE dataset_catalog (
  dataset_id TEXT PRIMARY KEY, dataset_name TEXT NOT NULL,
  competition_name TEXT, competition_year INTEGER,
  topic TEXT NOT NULL,
  data_structure TEXT CHECK (data_structure IN ('Cross-section', 'Time-Series', 'Panel')),
  geometry TEXT CHECK (geometry IN ('None', 'Spatial Weights Matrix', 'Coordinates', 'Polygons')),
  response_type TEXT CHECK (response_type IN ('Binary', 'Multiclass', 'Ordered', 'Count', 'Survival', 'Regression')),
  rows BIGINT CHECK (rows > 0), features INTEGER CHECK (features >= 0), size_mb DOUBLE CHECK (size_mb >= 0),
  licence TEXT, delivery TEXT CHECK (delivery IN ('Built-in', 'Download')), source_url TEXT,
  CHECK ((competition_name IS NULL) = (competition_year IS NULL)));

CREATE TABLE dataset_references (
  dataset_id TEXT REFERENCES dataset_catalog, bibtex_key TEXT,
  role TEXT CHECK (role IN ('Primary', 'Competition', 'Benchmark', 'Other')),
  PRIMARY KEY (dataset_id, bibtex_key));

CREATE TABLE model_catalog (
  model_id TEXT PRIMARY KEY, model_name TEXT NOT NULL,
  estimator_family TEXT NOT NULL, estimator_method TEXT NOT NULL,   -- a method of its family
  package TEXT, language TEXT CHECK (language IN ('Julia', 'Python', 'R', 'C/C++', 'Fortran')));

CREATE TABLE model_references (
  model_id TEXT REFERENCES model_catalog, bibtex_key TEXT,
  role TEXT CHECK (role IN ('Methodology', 'Software', 'Other')),
  PRIMARY KEY (model_id, bibtex_key));

CREATE TABLE benchmarks (
  run_id TEXT, model_id TEXT REFERENCES model_catalog, dataset_id TEXT REFERENCES dataset_catalog,
  split_protocol TEXT, metric TEXT, score DOUBLE,
  PRIMARY KEY (run_id, metric));                 -- one row per run and metric

CREATE VIEW model_comparisons AS
SELECT b.run_id, b.model_id, m.model_name, m.estimator_family, m.estimator_method, m.package, m.language,
  b.dataset_id, d.dataset_name, d.data_structure, d.geometry, d.response_type,
  b.split_protocol, b.metric, b.score
FROM benchmarks b JOIN model_catalog m USING (model_id) JOIN dataset_catalog d USING (dataset_id);
\end{lstlisting}

\subsection{A schema for each dataset}

For each dataset of the catalogue, \file{connect_db} looks for its folder by the rule of \S\ref{sec:folders}. When the
folder has Parquet files, it creates a schema named after the dataset id, with a view named after each file:

\begin{lstlisting}[language=SQL]
CREATE SCHEMA "tourism_2011";
CREATE VIEW "tourism_2011"."series"       AS SELECT * FROM read_parquet('.../tourism_2011/series.parquet');
CREATE VIEW "tourism_2011"."observations" AS SELECT * FROM read_parquet('.../tourism_2011/observations.parquet');
\end{lstlisting}

An id that is already the name of a schema of DuckDB stops the build with an error, rather than shadowing the schema.

\begin{center}\footnotesize
\renewcommand{\arraystretch}{1.2}
\begin{tabularx}{\linewidth}{@{}p{0.2\linewidth}L@{}}
\toprule
\textbf{View} & \textbf{Columns}\\
\midrule
\file{series} & \file{series} (its number), \file{competition} (the dataset id), \file{name} (as the source names it), \file{frequency}, \file{period} (1 for none), \file{horizon}, \file{domain}, \file{start}\\
\file{observations} & \file{series}, \file{t} (from 1), \file{value}, \file{split}: \file{train}, a value the forecasters saw, or \file{test}, one their forecasts are scored on\\
\file{data} & A row per case of a tabular dataset: the response, named by \file{target} in \file{info.toml}, and its features\\
\file{splits} & \file{row} (of \file{data}, from 1), \file{repeat}, \file{fold}: the fold whose test set the row is in\\
\bottomrule
\end{tabularx}
\end{center}

For example, the best tabular models on one dataset:

\begin{lstlisting}[language=SQL]
SELECT model_name, estimator_family, score
FROM model_comparisons
WHERE dataset_id = 'heloc' AND metric = 'roc_auc'
ORDER BY score DESC LIMIT 10;
\end{lstlisting}

\section{The folders on your computer}
\label{sec:folders}

\subsection{Where the data of a dataset are}

There is one rule for every dataset. Its folder is \file{data/builtin/<id>} of the package if that exists, else
\file{~/.pollis/competitions/<id>} when an \file{info.toml} is there. The data of a dataset whose delivery is Download go
into the second folder. The three pages only read these folders.

\begin{lstlisting}
<package>/data/builtin/<id>/       or    ~/.pollis/competitions/<id>/
  info.toml
  series.parquet                          a dataset of series (format monash)
  observations.parquet
  data.parquet                            a tabular dataset (format tabular)
  splits.parquet
\end{lstlisting}

\subsection{The runs of Score Models}

Each run is a folder of \file{~/.pollis/competitions/runs}, named after the dataset, the frequency and the time; a run of several datasets is named
\file{datasets-} and the time, such as \file{datasets-20261120-101500}.

\begin{lstlisting}
~/.pollis/competitions/runs/<run>/
  run.toml          the choices of step 1
  method.jl         your forecast(y, h, period), as step 2 held it
  run.jl            the script the page writes and the Julia REPL runs
  progress.json     {"done", "total", "seconds", "failed"}, every hundredth of the series
  forecasts.csv     series,h,value
  results.json      what step 4 draws, or {"error": ...} when the run stopped
  <step>.jl, <step>.json   the code and the result of each Next Step run on it
  <id>/             in a run of several datasets: a run of one dataset, as above, for each
\end{lstlisting}

\begin{lstlisting}
# run.toml of a run of one dataset
competition = "tourism_2011"        # the dataset id
folder = "..."                      # only for a dataset of your own: read in place of the id
frequency = "Monthly"               # "" for all its frequencies
sample = 0                          # 0 for all the series; else the number drawn
seed = 1
started = "2026-11-20 10:15"

# run.toml of a run of several datasets
datasets = ["tourism_2011", "nn5"]
sample = 0
seed = 1
started = "2026-11-20 10:15"
\end{lstlisting}

\file{results.json} holds \file{finished}, \file{series}, \file{whole}, \file{level} (the frequency, or
\file{Total}), \file{measures} (sMAPE, MASE, OWA), \file{method} (a score per measure) and \file{benchmarks} (a
\file{name} and \file{scores} each). It also holds \file{plot}, with the series, its last \file{past} values,
\file{actual} and \file{forecast}, and the run's \file{seconds}.

\subsection{Your datasets: data/forecasting of the Pollis Workspace}
\label{sec:workspace}

Manage Competition Datasets writes into the data folder of your Pollis Workspace and nowhere else:

\begin{lstlisting}
<Pollis Workspace>/data/forecasting/
  forecastingCatalogue.toml     the whole catalogue as the page last read it, for Julia; rewritten, never edited
  <id>/                         a dataset saved from the catalogue
    the files, copied
    saved.toml                  the size, time and SHA-256 of each file when saved, so a change is seen
  <id>/                         a dataset you added from a file of your own
    series.parquet
    observations.parquet
    dataset.toml                how it was converted, and its entry in the catalogue
    saved.toml
\end{lstlisting}

\begin{lstlisting}
# dataset.toml of a dataset you added
[conversion]                # change it with Edit Facts on the page, which converts the file again
file = "sales.csv"
series_column = "store"
time_column = "week"
value_column = "units"
frequency = "Weekly"
period = 52
horizon = 13                # the last 13 values of each series are held out
series = 45

[[datasets]]                # its entry in the catalogue, under the same schema as every other
dataset_id = "store_sales"
dataset_name = "Store Sales"
topic = "Retail"

  [[datasets.references]]
  bibtex_key = "..."
  role = "Primary"
\end{lstlisting}

\begin{lstlisting}
# saved.toml
saved = "2026-11-20T10:15:00.000Z"

[[files]]
path = "series.parquet"
bytes = 4096
mtime = 1763633700000
sha256 = "..."
\end{lstlisting}

Score Models lists a dataset you added after those of the package, and scores it the same way. The run's
\file{run.toml} gives its \file{folder}, and \file{PC.folder} returns that folder as it is.

\subsection{The studies of Agents' Forecasting Competition}

\begin{lstlisting}
<Pollis Workspace>/orchestration/forecasting-competition/<study>/
  study.toml        the dataset, frequency, sample and seed, and the agents
  prepared.toml     what Prepare the Series wrote
  data/
    series.csv      series,period,horizon,t,value: only the values the forecasters saw
    README.md       what the columns are
  agent-1/ ... agent-N/   2 to 8 agents, each with
    SKILL.md        the skill forecast-independently
    TASK.md         the series, the horizon, the rules, how to run Julia
    data/           a copy of the study's data/
    check.jl        written by Pollis: runs forecast on every series as the scoring will
    method.jl       the agent's forecast(y, h, period)
    REPORT.md       what the method does and why
  score.jl          written by Score the Methods
  results.json      {"finished", "series", "whole", "rows"}
\end{lstlisting}

\begin{lstlisting}
# study.toml
competition = "tourism_2011"
frequency = "Quarterly"
sample = 300                # 0 for all the series of the frequency
seed = 1
created = "2026-11-20 10:15"

[[agents]]
agent = "claude"
model = "..."
label = "..."

# prepared.toml
series = 300
horizon = 8
numbers = [3, 17, 21, ...]  # the series drawn, by their numbers in the dataset
\end{lstlisting}

The held-out values are never written into the study. Score the Methods reads them from the package after every agent
has ended. Each row of \file{results.json} is an agent, a combination (the mean or the median of the agents'
forecasts) or a benchmark, with \file{kind}, \file{name} and \file{scores}.

\section{Where to change what}

\begin{center}\footnotesize
\renewcommand{\arraystretch}{1.2}
\begin{tabularx}{\linewidth}{@{}p{0.44\linewidth}L@{}}
\toprule
\textbf{File} & \textbf{What}\\
\midrule
\file{PollisCompetitions.jl/metadata/*.toml} & The catalogues and the published scores\\
\file{PollisCompetitions.jl/src/} & Reading, scoring, the benchmarks and the database\\
\file{.../compose/browser/media/forecastingCatalogue.toml} & The schema of the catalogue: its fields and views\\
\file{.../compose/browser/forecastingCatalogue.ts} & Reading and checking the catalogue's contributions\\
\file{.../compose/browser/forecastingCataloguePage.ts} & Manage Competition Datasets\\
\file{.../compose/browser/forecastingWorkspace.ts} & \file{data/forecasting/}: saved and added datasets\\
\file{.../compose/browser/forecastingActive.ts} & The datasets ticked Active\\
\file{.../compose/browser/forecastingCompetitionsPage.ts} & Score Models: the runs and their scripts\\
\file{.../compose/browser/agentForecastingPage.ts} & Agents' Forecasting Competition: the studies and their scripts\\
\file{.../model/browser/webviews/forecasting*.toml}, \file{agentForecasting.toml} & The panel texts, wikis and references of the three pages\\
\file{pollis-agents/skills/forecast-independently/SKILL.md} & The rules the agents follow\\
\bottomrule
\end{tabularx}
\end{center}

\references
\dref{AN2000a}{Assimakopoulos, V. and Nikolopoulos, K. (2000). The Theta Model: A Decomposition Approach to Forecasting. \emph{International Journal of Forecasting} 16(4), 521--530. \ddoi{10.1016/s0169-2070(00)00066-2}}
\dref{GBWHM2021a}{Godahewa, R., Bergmeir, C., Webb, G. I., Hyndman, R. J. and Montero-Manso, P. (2021). Monash Time Series Forecasting Archive. arXiv. \ddoi{10.48550/arXiv.2105.06643}}
\dref{HK2006a}{Hyndman, R. J. and Koehler, A. B. (2006). Another Look at Measures of Forecast Accuracy. \emph{International Journal of Forecasting} 22(4), 679--688. \ddoi{10.1016/j.ijforecast.2006.03.001}}
\dref{MSA2020a}{Makridakis, S., Spiliotis, E. and Assimakopoulos, V. (2020). The M4 Competition: 100,000 Time Series and 61 Forecasting Methods. \emph{International Journal of Forecasting} 36(1), 54--74. \ddoi{10.1016/j.ijforecast.2019.04.014}}
\dref{RM2019a}{Raasveldt, M. and Mühleisen, H. (2019). DuckDB: an Embeddable Analytical Database. In \emph{Proceedings of the 2019 International Conference on Management of Data (SIGMOD)}, 1981--1984. \ddoi{10.1145/3299869.3320212}}

\end{document}
